\chapter{Simplest Formulation of the Lattice Action}

\section{Gauge Action}

The gauge action can be expressed in terms of closed loops. I will outline
the construction for the simpler case of the abelian $U(1)$ model in which
the link variables are commuting complex numbers instead of SU(3) matrices.
Consider the simplest Wilson loop, the $1\times1$ plaquette
$W_{\mu\nu}^{1\times1}$:
\begin{equation}
\label{eq:plaqexpA}
\begin{aligned}
W_{\mu\nu}^{1\times1} &= U_\mu(x)\, U_\nu(x+\hat{\mu})\, U_\mu^\dagger(x+\hat{\nu})\, U_\nu^\dagger(x) \\
&= e^{iag\big[ A_\mu(x+\hat{\mu}/2) + A_\nu(x+\hat{\mu}+\hat{\nu}/2) -
          A_\mu(x+\hat{\nu}+\hat{\mu}/2) - A_\nu(x+\hat{\nu}/2) \big]}
\end{aligned}
\end{equation}
Expanding about $x + (\hat{\mu}+\hat{\nu})/2$ gives
\begin{equation}
\label{eq:plaqexpB}
\begin{aligned}
&\exp\Big[ ia^2g(\partial_\mu A_\nu - \partial_\nu A_\mu) +
       \frac{ia^4g}{12}(\partial_\mu^3 A_\nu - \partial_\nu^3 A_\mu) + \ldots \Big] \\
&= 1 + ia^2gF_{\mu\nu} - \frac{a^4g^2}{2}F_{\mu\nu}F^{\mu\nu} + O(a^6) + \ldots
\end{aligned}
\end{equation}
The real and imaginary parts of the loop give
\begin{equation}
\label{eq:plaqexpC}
\begin{aligned}
\operatorname{Re}\operatorname{Tr}(1-W_{\mu\nu}^{1\times1}) &=
   \frac{a^4g^2}{2}F_{\mu\nu}F^{\mu\nu} + \text{terms higher order in } a \\
\operatorname{Im}(W_{\mu\nu}^{1\times1}) &= a^2gF_{\mu\nu} .
\end{aligned}
\end{equation}
So far the indices $\mu$ and $\nu$ are uncontracted. There are $6$
distinct positively oriented plaquettes, $\{\mu<\nu\}$, associated with
each site. Summing over $\mu,\nu$, and taking care of the double counting
by an extra factor of $\frac{1}{2}$, gives
\[
\frac{1}{g^2}\sum_x\sum_{\mu<\nu}\operatorname{Re}\operatorname{Tr}(1-W_{\mu\nu}^{1\times1})
= \frac{a^4}{4}\sum_x\sum_{\mu,\nu}F_{\mu\nu}F^{\mu\nu}
\to \frac{1}{4}\int d^4x\, F_{\mu\nu}F^{\mu\nu} .
\]
Thus, to lowest order in $a$, the expansion of the plaquette gives the
continuum action. The steps in the derivation for non-abelian groups are
identical and the final expression (up to numeric factors) is the same. The
result for the gauge action for SU(3), defined in terms of plaquettes and
called the Wilson action, is
\begin{equation}
\label{eq:Wgaugeaction}
S_g = \frac{6}{g^2}\sum_x\sum_{\mu<\nu}\operatorname{Re}\operatorname{Tr}\,
      \frac{1}{3}(1-W_{\mu\nu}^{1\times1}) .
\end{equation}
For historic reasons the lattice calculations are mostly presented in terms
of the coupling $\beta \equiv 6/g^2$. Since $\beta$ is used for many other
quantities like $1/kT$ or the $\beta$-function, I shall try to be more
explicit in the use of notation ($g$ versus $\beta$) if there is a possibility
for confusion.

\medskip
\noindent\textit{Exercise: Repeat the Taylor expansion for the physical case
of SU(3). Show that the non-abelian term in $F_{\mu\nu} = \partial_\mu A_\nu -
\partial_\nu A_\mu + g[A_\mu, A_\nu]$ arises due to the non-commuting property
of $\lambda$ matrices. The details of the derivation for SU(2) are given in
section 3 of Ref.~[8].}
\medskip

There are four important points to note based on the above construction of
the lattice action.

\textbf{1)} The leading correction is $O(a^2)$. The term
$\frac{a^2}{6}F_{\mu\nu}(\partial_\mu^3 A_\nu - \partial_\nu^3 A_\mu)$ is
present in the expansion of all planar Wilson loops. Thus at the classical
level it can be gotten rid of by choosing an action that is a linear
combination of say $1\times1$ and $1\times2$ Wilson loops with the appropriate
relative strength given by the Taylor expansion (see Section 12).

\textbf{2)} Quantum effects will give rise to corrections, \textit{i.e.}
$a^2 \to X(g^2)a^2$ where in perturbation theory
$X(g^2) = 1 + c_1 g^2 + \ldots$, and will bring in additional non-planar
loops. Improvement of the action will consequently require including these
additional loops, and adjusting the relative strengths which become
functions of $g^2$. This is also discussed in Section 12.

\textbf{3)} The reason for defining the action in terms of small loops is
computational speed and reducing the size of the discretization errors. For
example the leading correction to $1\times1$ loops is proportional to
$a^2/6$ whereas for $1\times2$ loops it increases to $5a^2/12$. Also, the
cost of simulation increases by a factor of $2-3$.

\textbf{4)} The electric and magnetic fields $E$ and $B$ are proportional
to $F_{\mu\nu}$. Eq.~\ref{eq:plaqexpC} shows that these are given in terms
of the imaginary part of Wilson loops.

\section{Fermion Action}

To discretize the Dirac action, Wilson replaced the derivative with the
symmetrized difference and included appropriate gauge links to maintain
gauge invariance
\begin{equation}
\label{eq:Dslnaive}
\bar\psi\not{D}\psi = \frac{1}{2a}\bar\psi(x)\sum_\mu\gamma_\mu
\big[ U_\mu(x)\psi(x+\hat{\mu}) - U_\mu^\dagger(x-\hat{\mu})
      \psi(x-\hat{\mu}) \big] .
\end{equation}
It is easy to see that one recovers the Dirac action in the limit $a \to 0$
by Taylor expanding the $U_\mu$ and $\psi(x+\hat{\mu})$ in powers of the
lattice spacing $a$. Keeping only the leading term in $a$,
Eq.~\ref{eq:Dslnaive} becomes
\[
\begin{aligned}
&\frac{1}{2a}\bar\psi(x)\gamma_\mu\Big[
   \big(1+iagA_\mu(x+\hat{\mu}/2)+\ldots\big)
   \big(\psi(x)+a\psi'(x)+\ldots\big) \\
&\quad - \big(1-iagA_\mu(x-\hat{\mu}/2)+\ldots\big)
   \big(\psi(x)-a\psi'(x)+\ldots\big) \Big] \\
&= \bar\psi(x)\gamma_\mu\Big(\partial_\mu + \frac{a^2}{6}\partial_\mu^3 +
   \ldots\Big)\psi(x) \\
&\quad + ig\bar\psi(x)\gamma_\mu\Big[A_\mu + \frac{a^2}{2}\Big(
   \frac{1}{4}\partial_\mu^2 A_\mu + (\partial_\mu A_\mu)\partial_\mu +
   A_\mu\partial_\mu^2\Big) + \ldots \Big]\psi(x) ,
\end{aligned}
\]
which, to $O(a^2)$, is the kinetic part of the standard continuum Dirac
action in Euclidean space-time. Thus one arrives at the simplest (called
``naive'') lattice action for fermions
\begin{equation}
\label{eq:naiveact}
\begin{aligned}
\mathcal{S}^N &= m_q\sum_x\bar\psi(x)\psi(x) \\
&\quad + \frac{1}{2a}\sum_x\bar\psi(x)\gamma_\mu\big[
   U_\mu(x)\psi(x+\hat{\mu}) - U_\mu^\dagger(x-\hat{\mu})
   \psi(x-\hat{\mu}) \big] \\
&\equiv \sum_x\bar\psi(x)M_{xy}^N[U]\psi(y)
\end{aligned}
\end{equation}
where the interaction matrix $M^N$ is
\begin{equation}
\label{eq:naiveactM}
M_{i,j}^N[U] = m_q\delta_{ij} + \frac{1}{2a}\sum_\mu
\big[ \gamma_\mu U_{i,\mu}\delta_{i,j-\mu} -
      \gamma_\mu U_{i-\mu,\mu}^\dagger\delta_{i,j+\mu} \big]
\end{equation}
The Euclidean $\gamma$ matrices are hermitian, $\gamma_\mu = \gamma_\mu^\dagger$,
and satisfy $\{\gamma_\mu,\gamma_\nu\} = 2\delta_{\mu\nu}$. The
representation I shall use is
\begin{equation}
\label{eq:gammamatdef}
\vec\gamma = \begin{pmatrix} 0 & i\vec\sigma \\ -i\vec\sigma & 0 \end{pmatrix},
\qquad
\gamma_4 = \begin{pmatrix} 1 & 0 \\ 0 & -1 \end{pmatrix},
\qquad
\gamma_5 = \begin{pmatrix} 0 & 1 \\ 1 & 0 \end{pmatrix} ,
\end{equation}
which is related to Bjorken and Drell conventions as follows:
$\gamma_i = i\gamma_{BD}^i$, $\gamma_4 = \gamma_{BD}^0$,
$\gamma_5 = \gamma_{BD}^5$. In this representation $\gamma_1, \gamma_3$ are
pure imaginary, while $\gamma_2, \gamma_4, \gamma_5$ are real.

The Taylor expansion showed that the discretization errors start at
$O(a^2)$. For another simple illustration consider the inverse of the
free-field propagator $m + \frac{i}{a}\sum_\mu\gamma_\mu\sin(p_\mu a)$. Set
$\vec{p}=0$ and rotate to Minkowski space ($p_4 \to iE$, \textit{i.e.}
$\sin p_4 a \to i\sinh Ea$). Then, using the forward propagator (upper two
components of $\gamma_4$), gives
\begin{equation}
\label{eq:Npolemass}
m_q^{\mathrm{pole}} a = \sinh Ea
\end{equation}
for the relation between the pole mass and the energy. This shows that, even
in the free field case, the continuum relation $E(\vec{p}=0) = m$ is violated
by corrections that are $O(a^2)$.

\noindent\textbf{Symmetries}: The invariance group of the fermion action
under rotations in space and time is the hypercubic group. Full Euclidean
invariance will be recovered only in the continuum limit. The action is
invariant under translations by $a$ and under $\mathcal{P}$, $\mathcal{C}$,
and $\mathcal{T}$ as can be checked using Table 2.

The naive action $\bar\psi_x M_{xy}^N\psi_y$ has the following global
symmetry:
\begin{equation}
\label{eq:vector_psi}
\begin{aligned}
\psi(x) &\to e^{i\theta}\psi(x) \\
\bar\psi(x) &\to \bar\psi(x)e^{-i\theta}
\end{aligned}
\end{equation}
where $\theta$ is a continuous parameter. This symmetry is related to baryon
number conservation and leads to a conserved vector current. For $m_q = 0$
the action is also invariant under
\begin{equation}
\label{eq:axial_psi}
\begin{aligned}
\psi(x) &\to e^{i\theta\gamma_5}\psi(x) \\
\bar\psi(x) &\to \bar\psi(x)e^{i\theta\gamma_5}
\end{aligned}
\end{equation}
Having both a vector and an axial symmetry in a hard-cutoff regularization
would imply a violation of the Adler-Bell-Jackiw theorem. It turns out that
while the naive fermion action preserves chiral symmetry it also has the
notorious fermion doubling problem as explained in Section 9. The chiral
charges of these extra fermions are such as to exactly cancel the anomaly.
In fact the analysis of $\mathcal{S}^N$ lead to a no-go theorem by
Nielsen-Ninomiya [14] that states that it is not possible to define a local,
translationally invariant, hermitian lattice action that preserves chiral
symmetry and does not have doublers.

We will discuss two fixes to the doubling/anomaly/chiral symmetry problem.
First, include an additional term in the action that breaks chiral symmetry
and removes the doublers (Wilson's fix). Second, exploit the fact that the
naive fermion action has a much larger symmetry group, $U_V(4)\otimes U_A(4)$,
to reduce the doubling problem from $2^d = 16 \to 16/4$ while maintaining a
remnant chiral symmetry (staggered fermions). In both cases one regains the
correct anomaly in the continuum limit. Before diving into these fixes, I
would like to first discuss issues of the measure of integration and gauge
fixing, and give an overview of Symanzik's improvement program.

\section{The Haar Measure}

The fourth ingredient needed to complete the formulation of LQCD as a quantum
field theory via the path integral is to define the measure of integration
over the gauge degrees of freedom. Note that, unlike the continuum fields
$A_\mu$, lattice fields are SU(3) matrices with elements that are bounded in
the range $[0,1]$. Therefore, Wilson proposed an invariant group measure, the
Haar measure, for this integration. This measure is defined such that for any
elements $V$ and $W$ of the group
\begin{equation}
\label{eq:haarmeasure}
\int dU\, f(U) = \int dU\, f(UV) = \int dU\, f(WU)
\end{equation}
where $f(U)$ is an arbitrary function over the group. This construction is
simple, and has the additional advantage that in non-perturbative studies it
avoids the problem of having to include a gauge fixing term in the path
integral. This is because the field variables are compact. Hence there are no
divergences and one can normalize the measure by defining
\begin{equation}
\label{eq:haarnorm}
\int dU = 1 .
\end{equation}
For gauge invariant quantities each gauge copy contributes the same amount.
Consequently the lack of a gauge fixing term just gives an overall extra
numerical factor. This cancels in the calculation of correlation functions
due to the normalization of the path integral. Note that gauge fixing and the
Fadeev-Papov construction does have to be introduced when doing lattice
perturbation theory [36].

The measure and the action are invariant under $\mathcal{C}$, $\mathcal{P}$,
$\mathcal{T}$ transformations, therefore when calculating correlation
functions via the path integral we must average over a configuration
$\mathcal{U}$ and the ones obtained from it by the application of
$\mathcal{C}$, $\mathcal{P}$, $\mathcal{T}$. In most cases this averaging
occurs very naturally, for example, Euclidean translation invariance
guarantees $\mathcal{P}$ and $\mathcal{T}$ invariance. Similarly, taking the
real part of the trace of a Wilson loop is equivalent to summing over $U$ and
its the charge conjugate configuration. A cook-book recipe for evaluating the
behavior of correlation functions after average over configurations related
by the lattice discrete symmetries is given in Section 17.

One important consequence of these symmetries is that only the real part of
Euclidean correlation functions has a non-zero expectation value. This places
restriction on the quantities one can calculate using LQCD. For example,
simulations of LQCD have not been very successful at calculations of
scattering amplitudes, which are in general complex [48]. The connection
between Minkowski and Euclidean correlation functions was discussed in
Section 5. To get scattering amplitudes what is required is to first
calculate Euclidean amplitudes for all values of the Euclidean time $\tau$,
followed by a fourier transform in $\tau$, and finally a Wick rotation to
Minkowski space. For such an procedure to work requires data with phenomenal
accuracy. Unfortunately, simulations give Euclidean correlation functions for
a discrete set of $\tau$ and with errors that get exponentially magnified in
the rotation. The best that has been done is to use L\"uscher's method that
relates scattering phase shifts to shifts in two particle energy states in
finite volume [27,28]. So far this method has been tested on models in 2--4
dimensions, and used to calculate scattering lengths in QCD [29,30].

\section{Gauge Fixing and Gribov Copies}

Gauge invariance and the property of group integration, $\int dU\, U = 0$,
are sufficient to show that only gauge invariant correlation functions have
non-zero expectation values. This is the celebrated Elitzur's theorem [37].
It states that a local continuous symmetry cannot be spontaneously broken.

In certain applications (study of quark and gluon propagators [38], 3-gluon
coupling [39], hadronic wavefunctions [40], and determination of
renormalization constants using quark states [41]) it is useful to work in a
fixed gauge. The common gauge choices are axial ($\eta_\mu A_\mu = 0$ where
$\eta$ in Minkowski space is a fixed time-like vector. The common choice is
$A_4 = 0$), Coulomb ($\partial_i A_i = 0$), and Landau
($\partial_\mu A_\mu = 0$). Their lattice analogues, assuming periodic
boundary conditions, are the following.

\textbf{Maximal Axial Gauge}: On a lattice with periodic boundary conditions
one cannot enforce $A_4 = 0$ by setting all links in the time direction to
unity as the Wilson line in each direction is an invariant. The gauge freedom
at all but one site can be removed as follows. Set all time-like links $U_4 =
1$ on all but one timeslice. On that time-slice all $U_4$ links are
constrained to be the original Wilson line in 4-direction, however, one can
set links in $\hat{z}$ direction $U_3 = 1$ except on one z-plane. On that
z-plane set $U_3$ links to the Wilson line in that direction and all links in
$\hat{y}$ direction $U_2 = 1$ except on one y-line. On this y-line set all
$U_2$ links to the value of y-lines and links in $\hat{x}$ direction $U_1 = 1$
except on one point. While such a construction is unique, the choice of the
time-slice, z-plane, y-line, and x-point and the labeling of $x,y,z,t$ axis
are not. Also, this gauge fixing does not respect Euclidean translation
invariance and it is not smooth.

\textbf{Coulomb Gauge}: The lattice condition is given by the gauge
transformation $V(x)$ that maximizes the function
\begin{equation}
\label{eq:Coulombgauge}
F[V] = \sum_x\sum_i \operatorname{Re}\operatorname{Tr} V(x)
\bigg( U_i(x)V^\dagger(x+\hat{i}) +
       U_i^\dagger(x-\hat{i})V^\dagger(x-\hat{i}) \bigg)
\end{equation}
separately on each time-slice and $i$ runs over the spatial indices $1-3$.

\textbf{Landau Gauge}: This is defined by the same function as in
Eq.~\ref{eq:Coulombgauge} except that the sum is evaluated on the whole
lattice and $i = 1-4$.

The maximization condition, Eq.~\ref{eq:Coulombgauge}, corresponds, in the
continuum, to finding the stationary points of $F[V] = \lVert A^V \rVert =
\int d^4x\, \operatorname{Tr}(A_\mu^V)^2(x)$. What we would like is to find
the global minimum corresponding to the ``smoothest'' fields $A_\mu$, or the
copies in case it is not unique. For any gauge configuration there is an
infinite set of gauge equivalent configurations obtained by applying the set
of transformations $\{V(x)\}$. This set is called the gauge orbit. The first
question is whether the gauge fixing procedure gives a unique solution,
\textit{i.e.} whether the gauge-fixing algorithm converges to the same final
point irrespective of the starting point on the gauge orbit. For the maximal
axial gauge, the construction is unique (the configuration $\{U\}$ and
$\{VUV^\dagger\}$ for any $\{V\}$ give the same final link matrices), so the
answer is clearly yes. For the Coulomb and Landau gauge choices the answer is
no --- there are a number of solutions to the gauge-fixing condition. These
are the famous Gribov copies [42]. For a recent theoretical review of this
subject see [43].

The algorithms for fixing to either the Coulomb gauge or the Landau gauge are
mostly local. Local algorithms generate, for each site, a gauge matrix
$v_i(x)$ that maximizes the sum, defined in Eq.~\ref{eq:Coulombgauge}, of the
link matrices emerging from that site. This gauge transformation is then
applied to the lattice. These two steps are iterated until the global maximum
is found. The ordered product $\prod_i v_i(x)$ gives $V(x)$. For all
practical purposes we shall call this end-point of the gauge fixing algorithm
a Gribov copy. It is, in most cases and especially for the local gauge-fixing
algorithms, the nearest extremum, and not the global minimum of
$\lVert A^V \rVert$. A set of Gribov copies can be generated by first
performing a different random gauge transformation on the configuration and
then applying the same gauge-fixing algorithm, or by using different
gauge-fixing algorithms. For recent studies of the efficacy of different
algorithms see [44,45].

The Monte Carlo update procedure is such that it is highly improbable that
any finite statistical sample contains gauge equivalent configurations. So
one does not worry about gauge copies. In the gauge-fixed path integral one
would \textit{a priori} like to average over only one Gribov copy per
configuration and choose this to be the ``smoothest''. Finding the $V(x)$
which gives the smoothest configuration has proven to be a non-trivial
problem.

The open questions are: (i) how to define the ``smoothest'' solution of the
gauge-fixing algorithm (\textit{i.e.} defining the fundamental modular
domain), (ii) how to ascertain that a solution is unique (lies within or on
the boundary of the fundamental modular domain), and (iii) how to find it.

The ultimate question is --- what is the dependence of the quantities one
wants to calculate on the choice of Gribov copy? A number of studies have
been carried out [46], however the results and their interpretation are not
conclusive. Firstly, it is not known whether the smoothest solution lying
within the fundamental modular domain can be found economically for all
configurations and over the range of gauge coupling being investigated. In
the absence of the smoothest solution it is not clear whether (i) the
residual gauge freedom simply contributes to noise in the correlation
function, and (ii) one should average over a set of Gribov copies for each
configuration to reduce this noise. In my opinion these remain challenging
problems that some of you may wish to pursue.

The problem of Gribov copies and the Faddeev-Papov construction (a gauge
fixing term and the Faddeev-Papov determinant) resurfaces when one wants to
do lattice perturbation theory or when taking the Hamiltonian limit of the
path integral via the transfer matrix formalism. This construction is
analogous to that in the continuum [36]. Finally, let me mention that in a
fixed gauge, the lattice theory has only a BRST invariance [27]. While BRST
invariance is sufficient to prove that the theory is renormalizable, the set
of operators of any given dimension allowed by BRST invariance is larger than
that with full gauge invariance. In the analyses of improvement of the action
and operators it is this enlarged set that has to be considered [47].
